\subsubsection{Examples of Visual Recognition Tasks}\label{sec:examples-of-visual-recognition-tasks}

Now that we know what Computer Vision is, the next question is: \emph{\textbf{what kinds of problems does Computer Vision solve?}} Here, we present a few examples of visual recognition tasks that are commonly solved using Computer Vision techniques. These tasks differ mainly in \textbf{what information the algorithm is expected to produce}.
\begin{itemize}
    \item \important{Object Detection}. Given an image $I$ find \textbf{what} objects exist and \textbf{where} they are located. Mathematically:
    \begin{equation*}
        I \to \left\{
            \left(c_i , b_i\right)
        \right\}_{i=1}^{N}
    \end{equation*}
    Where $c_i$ is the class of the $i$-th object and $b_i$ is the bounding box of the $i$-th object.

    For example, given an image of a street scene, the algorithm should be able to identify and locate cars, pedestrians, traffic lights, etc. Unlike image classification, there may be multiple objects, classes and instances of the same class in the image (i.e., multiple cars, pedestrians, etc.).

    \textcolor{Red2}{\faIcon{question-circle} \textbf{Why is difficult?}} The detector must answer simultaneously ``\emph{what is this?}'' and ``\emph{where is it?}''. It is therefore solving both \textbf{recognition} and \textbf{localization}.


    \item \important{Pose Estimation}. Instead of recognizing objects, the goal is to understand the \textbf{configuration of the human body}. Rather than returning a label ``\emph{Person}'' or a bounding box, the algorithm predicts the locations of key points on the human body, such as joints (e.g., elbows, knees, shoulders). Mathematically:
    \begin{equation*}
        I \to \left\{
            \left(x_i, y_i\right)
        \right\}_{i=1}^{K}
    \end{equation*}
    Where $K$ is the number of key points and $(x_i, y_i)$ are the coordinates of the $i$-th key point. Thus, the algorithm is \textbf{not classifying} a person, but rather it is estimating \textbf{their geometry}.


    \item \important{Quality Inspection}. Factories continuously produce components, some of which contain defects. The goal of quality inspection is to use a vision system to automatically identify normal products, defective products, and defect types instead of asking a human operator to inspect thousands of products.
    
    Here, the model can perform classification, and even \textbf{novel-class detection}. This means that if the defect has never been observed before, the system should recognize ``\emph{this does not belong to any known class}''. This capability is extremely valuable in manufacturing, where unexpected failure modes can appear.


    \item \important{Image Captioning}. Instead of producing a class, the model produces an \textbf{entire sentence}. For example, given an image of a dog playing with a ball, the model might generate the caption: ``\emph{A dog is playing fetch with a ball in the park.}''. Mathematically:
    \begin{equation*}
        I \to \emph{Sentence}
    \end{equation*}
    \textcolor{Red2}{\faIcon{question-circle} \textbf{Why is difficult?}} Now the model must combine Computer Vision with Natural Language Processing (NLP). It must understand objects, actions, relationships, grammar, word order. This is much harder than classification.
\end{itemize}
In general, even very advanced Computer Vision systems still make mistakes because visual understanding remains a \textbf{challenging problem}.

\highspace
\begin{flushleft}
    \textcolor{Green3}{\faIcon{check-circle} \textbf{The evolution of Computer Vision}}
\end{flushleft}
Years ago, algorithms were designed manually. Researchers created mathematical models describing edges, textures, corners, gradients and other features. The focus was on explicitly engineering features.

\highspace
Nowadays, Machine Learning, and especially Deep Learning, has become the dominant paradigm. Instead of manually defining useful features, the \textbf{model learns them directly from data}.
\begin{itemize}
    \item \textbf{Before}: mathematical/statistical descriptions of images.
    \item \textbf{Now}: machine-learning methods are much more popular.
\end{itemize}
\textcolor{Green3}{\faIcon{question-circle} \textbf{Thus, why start from Image Classification?}} If Computer Vision can do so many things, \emph{why do we start with Image Classification?} The answer is that Image Classification is the \textbf{simplest visual recognition problem}. The input is only an image, and the output is only a class. Although it is simple, it is still a very useful task. Almost every advanced Computer Vision task extends this basic problem.